\documentclass[11pt]{article}
\usepackage[a4paper,margin=22mm]{geometry}
\usepackage{amsmath,amssymb,booktabs,array,tabularx,graphicx,xcolor,hyperref}
\hypersetup{colorlinks=true,linkcolor=blue!55!black,urlcolor=blue!55!black}
\setlength{\parindent}{0pt}
\setlength{\parskip}{5pt}
\setlength{\emergencystretch}{2em}
\newcommand{\code}[1]{\texttt{#1}}
\title{Evidence-first PRL LangGraph\\Implementation and Readiness Report}
\author{Active Learning FFLO / Topology Project}
\date{2026-08-07}
\begin{document}
\maketitle

\begin{center}
\fcolorbox{blue!55!black}{blue!4}{\parbox{0.92\textwidth}{\textbf{Status note
(2026-08-12).} This document records the original architecture baseline. Its
Stage I--III recommendation and pending scope/claim status are superseded by
\code{audits/MANUSCRIPT\_SCOPE.md} and the dated claim-approval record. The
current approved scope spans Stages I--IV and the workflow stops at
\code{GATE\_FIGURE\_APPROVAL}.}}
\end{center}

\begin{abstract}
An isolated evidence-first manuscript workflow has been built and revalidated
for the Active Learning FFLO project. The package preserves the supplied
prompt, audits the available artifacts, proposes manuscript scope and figures,
freezes claims through explicit registries, registers a 15-node LangGraph, and
retains a Stage III Results-by-figure vertical slice as an approval package. It
does not modify physics code, physical definitions, frozen data, or QAHE3D. The
default run stops at \code{GATE\_SCOPE\_APPROVAL}.
\end{abstract}

\section{Repository audit}
The curated manifest contains 37 entries: 34 present and three explicitly
missing. It covers representative reports and audit PDFs, frozen datasets,
convergence records, phase/topology maps, plotting and report sources, \LaTeX{}
sources, AI-workflow documentation, the recorded prompt, and LangGraph source.
No pre-existing paper manuscript draft was found. The missing records are the
Stage II endpoint, Stage III endpoint, and Stage III convergence-report bundle;
none was reconstructed from prose.

\section{Scientific scope assessment}
The strongest current line is: \emph{Cold-start active learning reconstructs
thermodynamic and topological FFLO phase diagrams using a numerically audited
BdG oracle.} Active learning remains a scheduler for exact BdG calls. Stages
I--III form the main narrative: the two-dimensional discovery problem, exact
oracle reliability and thermodynamic closure, and the audited Stage III
cFFLO/tFFLO result.

Stage IV is conditional and may be used only as a visibly non-converged 3D
demonstration. Stage V belongs in Discussion because V3/V4 superiority is not
established. Stage VI belongs in Outlook because the world model is not ready
for EIG and the requested spectral descriptors are absent.

\section{Manuscript workflow design}
The graph uses evidence order rather than article reading order:
\begin{center}
\small\ttfamily
artifact/provenance $\rightarrow$ readiness $\rightarrow$ scope/claims\\
$\rightarrow$ figures/Results $\rightarrow$ Methods/Supplement/Discussion\\
$\rightarrow$ Introduction $\rightarrow$ Abstract/Title $\rightarrow$ audits
$\rightarrow$ LaTeX $\rightarrow$ final approval
\end{center}
The staged boundaries are \code{GATE\_SCOPE\_APPROVAL},
\code{GATE\_CLAIM\_APPROVAL}, \code{GATE\_FIGURE\_APPROVAL}, permission-gated
Introduction retrieval, and \code{GATE\_FINAL\_MANUSCRIPT\_APPROVAL}. Every
Results paragraph records observation, interpretation, evidence source, figure
reference, caveats, and unsupported-sentence flags.

\section{Publication-readiness gaps}
The local checkout does not provide the Stage II/III final raw endpoints or the
Stage III convergence/figure bundle. It also does not provide matched-budget
random/Sobol baselines, a multi-seed benchmark, or an independent
hidden-ground-truth evaluator.

The strongest available evidence remains the signed-off audit layer. Phase II
passed the publication-boundary audit. Stage III returned Decision A under the
configured audit, with final coverage $p_{95}=0.006185<0.00625$ and trusted
surprise below 0.02 over the final three transitions. A $k=6$ sensitivity gives
coverage $p_{95}=0.006269$, slightly above the threshold, while map and
component structure remain stable. This sensitivity must remain visible.

\section{Four-figure plan and regeneration needs}
\begin{table}[htbp]
\small
\begin{tabularx}{\textwidth}{p{0.07\textwidth}p{0.23\textwidth}X p{0.14\textwidth}}
\toprule
ID & Scientific role & Required evidence and caveat & Status \\
\midrule
Fig. 1 & Model, exact oracle, label hierarchy, and AL loop & Canonical model/numerics; schematic is not numerical evidence & Ready \\
Fig. 2 & 2D thermodynamic discovery and closure & Recover frozen endpoint; retain hard-risk layer & Blocked \\
Fig. 3 & Stage III cFFLO/tFFLO result and convergence & Recover final raw/figure bundle; state $k$-sensitivity & Blocked \\
Fig. 4 & Benchmark or caveated 3D extension & Prefer matched multi-seed benchmark; otherwise label Stage IV non-converged & Conditional \\
\bottomrule
\end{tabularx}
\end{table}

Figures 2 and 3 remain blocked pending frozen raw bundles. Figure 1 is ready
for a reproducible vector redraw. Figure 4 requires a scope decision and must
be omitted if neither a caveated Stage IV demonstration nor an authorized
matched benchmark is available. Final scientific figures cannot be manually
edited without their source scripts.

\section{Graph execution semantics}
Evidence and claims are frozen before Results. Results are drafted figure by
figure with paragraph-level trace metadata. Methods, Supplement, and Discussion
are constrained by approved Results. The Introduction citation loop cannot run
until Results claims and external retrieval are approved. Abstract/Title are
generated last. Conditional edges stop execution at unresolved scope, claim,
figure, retrieval, or final human boundaries.

\begin{figure}[ht]
\centering
\includegraphics[width=0.84\textwidth]{figures/workflow_overview.pdf}
\caption{Evidence-first PRL generation order. Red nodes preserve human or
permission boundaries. No node can promote a scientific claim or fabricate a
missing artifact.}
\end{figure}

\section{QAHE3D adaptation boundary}
Only reusable semantics were adapted: Introduction claim maps,
permission-gated retrieval, version/full-text resolution, claim--citation
mapping, placement audit, paragraph blueprints, and agent governance concepts.
File-level upstream SHA-256 hashes are recorded. No QAHE3D code is imported at
runtime and no QAHE3D file was changed. The Active Learning adapter requires
approved scope, claims, and figure-traced Results.

\section{Stage III vertical slice}
The vertical slice reaches \code{GATE\_FIGURE\_APPROVAL}. It contains a claim
contract, figure contract, evidence extract, caption, three Results paragraphs,
paragraph metadata, consistency stub, audit report, and an isolated \LaTeX{}
fragment. The metadata explicitly separates observation and interpretation and
records source and figure references. All unsupported-sentence lists are empty.
The fragment compiles and renders cleanly, but its figure remains a QA
placeholder because the frozen figure source is absent.

\section{Readiness decision and next steps}
The architecture is ready for scientific review, but the default invocation
stops at:
\begin{center}
\small\ttfamily GATE\_SCOPE\_APPROVAL
\end{center}
After scope approval, the graph stops independently at claim and figure
approval. First choose the physics-first Stages I--III scope, then review the
claim registry and Stage III slice. Recover frozen Stage II/III bundles without
recomputation. A methods-efficiency positioning additionally requires separately
authorized matched baselines, multiple seeds, and hidden ground truth.

\section{Validation}
The graph registers 15 nodes and emits 16 gate records on a fully traversed test
invocation; the unapproved invocation stops at scope. Fourteen automated tests
pass. The Stage III QA fragment compiles to two pages and was visually
inspected. Prompt integrity was reverified against the SHA-256 recorded in the
prompt provenance file. No exact calculation, external retrieval, production
manuscript merge, or physics-definition change was performed.

\end{document}
